% kuleuventheme2 by Janez Kren, September 2017, janez.kren@kuleuven.be, based on:
% kuleuventheme 1.3 by Roland Pastorino, 2013 roland.pastorino@kuleuven.be / www.rolandpastorino.com

\PassOptionsToPackage{usenames, table, dvipsnames}{xcolor} % 
\documentclass[11pt,t]{beamer}
\usetheme{kuleuven2}	%THEME OPTIONS for LOGO: kul (default), kulak, lrd,    ; OPTIONS for TITLE PAGE: normal (default), sedes


%%% OTHER SETTINGS
\usefonttheme[onlymath]{serif}			% math font with serifs, delete to make it sans-serif
\setbeamertemplate{footline}[body] 		% delete this line to remove footline bar on all frames
%\usepackage[orientation=landscape,size=custom,width=16,height=9,scale=0.5,debug]{beamerposter} %enable for widescreen 16:9 ratio
%\titlegraphic{ \includegraphics[width=.2\paperwidth]{mytitlepagepic.png} } %optional title page image


%%% ADDED BY ALEX (start)
\usepackage[english]{babel}
\usepackage{amsfonts}
\usepackage{amssymb}

\usepackage{pythonhighlight}

\usepackage[orientation=landscape,size=custom, width=16,height=9,scale=0.5,debug]{beamerposter} 

\definecolor{mygray}{rgb}{0.8,0.8,0.8}
\newcommand{\inline}[1]{ \colorbox{mygray}{\lstinline{#1}}}

\usepackage{xcolor}

\usepackage{layouts} % To get text/page width

\usepackage{listings}
\definecolor{codegreen}{rgb}{0,0.6,0}
\definecolor{codegray}{rgb}{0.5,0.5,0.5}
\definecolor{codepurple}{rgb}{0.58,0,0.82}
% \definecolor{backcolour}{rgb}{0.95,0.95,0.92}
\definecolor{backcolour}{rgb}{0.93,0.93,0.95}
\lstdefinestyle{mystyle}{
    backgroundcolor=\color{backcolour},   
    commentstyle=\color{codegreen},
    keywordstyle=\color{magenta},
    numberstyle=\tiny\color{codegray},
    stringstyle=\color{codepurple},
    basicstyle=\ttfamily\footnotesize,
    breakatwhitespace=false,         
    breaklines=true,                 
    captionpos=b,                    
    keepspaces=true,                 
    numbers=left,                    
    numbersep=5pt,                  
    showspaces=false,                
    showstringspaces=false,
    showtabs=false,                  
    tabsize=2
}
\lstset{style=mystyle}

\usepackage{mathtools}
\DeclarePairedDelimiter\ceil{\lceil}{\rceil}

\usepackage{listings}

%%% ADDED BY ALEX (stop)


%%% TITLE PAGE INFO:
\title[Geavanceerde Computerarchitectuur 3$^{\text{rd}}$ session]{Geavanceerde Computerarchitectuur 3$^{\text{rd}}$ session} %[]] will appear in footline
\subtitle{Optimization techniques, \\ Partitioning the grid into warps and reducing thread divergence}

\author{Alex Marinsek, {\footnotesize alexander.marinsek@kueluven.be}}
\institute{DRAMCO -- WaveCore -- ESAT}
\date{Academic year 2023/24}



\begin{document}
\csname beamer@calculateheadfoot\endcsname %recalculate head and foot dimension


 %%
 %%  0. TITLE PAGE and TABLE OF CONTENT
 %%
% Title page
\begin{frame}[plain,noframenumbering]
	\titlepage
\end{frame}
	

% Table of Contents
\begin{frame}{Outline}
	\hfill	{\large \parbox{.961\textwidth}{\tableofcontents[hideothersubsections]}}
\end{frame}




\section{Background}


\begin{frame}{A warp}
    \begin{itemize}\setlength\itemsep{1em}
        \item A warp is a group of 32 threads
        \item $\text{number of warps} = \ceil*{ \frac{\text{Threads per block}}{\text{32}}} \cdot \text{number of blocks}$
        \item For example
        \begin{itemize}\setlength\itemsep{0.5em}\vspace{1mm}
            \item \inline{foo<<<2,64>>>()} \uncover<2->{$\Rightarrow$ 4 (128 threads \dots 128 cores)}
            \item \inline{foo<<<4,33>>>()} \uncover<3->{$\Rightarrow$ 8 (132 threads \dots \textbf{256 cores})}
        \end{itemize}
        \uncover<4->{\item Also called \textit{a sub-group} in OpenCL, a \textit{wavefront} by AMD, \dots}
    \end{itemize}
\end{frame}


\begin{frame}{Threads are grouped into warps and mapped to cores}
    \vspace{-2mm}
    
    \begin{figure}
        \centering
        \resizebox{0.9\textwidth}{!}{

\begin{tikzpicture}[]



%%% DRAWING %%%%%%%%%%%%%%%%%%%%%%%%%%%%%%%%%%%%%%%%%%%%%%%%%%%%%%%%%%%%%%%%%


%%% SW/FW and HW nodes

\node[] (grid) at (0, 0) {Grid};

\node[] (block_1) at (3, 0) {Block$_\text{1}$};
\node[] (block_2) at (3, -1) {Block$_\text{2}$};
\node[] at (3, -1.75) {\dots};
\node[] (block_n) at (3, -2.5) {Block$_\text{N}$};

\node[] (warp_1) at (6, 0) {Warp$_\text{1,1}$};
\node[] (warp_2) at (6, -1) {Warp$_\text{1,2}$};
\node[] at (6, -1.75) {\dots};
\node[] (warp_m) at (6, -2.5) {Warp$_\text{1,M}$};

\node[] (thread_1) at (9, 0) {Thread$_\text{1,1,1}$};
\node[] (thread_2) at (9, -1) {Thread$_\text{1,1,2}$};
\node[] at (9, -1.75) {\dots};
\node[] (thread_32) at (9, -2.5) {Thread$_\text{1,1,32}$};

% Two SMs
\foreach \x in {12, 14.75}{

    \draw[rounded corners, fill=black!20!white] (\x-1.2, -3.8) rectangle ++(2.4,5);
    \draw[rounded corners, fill=green] (\x-1, -3) rectangle ++(2,3.5);
    \draw[rounded corners, fill=green] (\x-1, -3.7) rectangle ++(2,0.6);
    
    \node[] at (\x, 0.8) {SM};
    \node[] at (\x, 0.8) {SM};
    \node[align=center] at (\x, -1.75) {32\\cores};
    \node[] at (\x, -3.5) {\dots};
}

%%% Lines (connections)

\draw[-{angle 60}] (grid.east) -- (block_1.west);
\draw[-{angle 60}, black!30!white] (grid.east) -- (block_2.west);
\draw[-{angle 60}, black!30!white] (grid.east) -- (block_n.north west);

\draw[-{angle 60}] (block_1.east) -- (warp_1.west);
\draw[-{angle 60}, black!30!white] (block_1.east) -- (warp_2.west);
\draw[-{angle 60}, black!30!white] (block_1.east) -- (warp_m.north west);

\draw[-{angle 60}] (warp_1.east) -- (thread_1.west);
\draw[-{angle 60}] (warp_1.east) -- (thread_2.west);
\draw[-{angle 60}] (warp_1.east) -- (thread_32.north west);

\draw[-{angle 60}] (thread_1.east) --++(1.3,0);
\draw[-{angle 60}] (thread_2.east) --++(1.3,0);
\draw[-{angle 60}] (thread_32.east) --++(1.2,0);

\draw[-{angle 60}, black!30!white] (warp_2.south east) to[out=-45, in=180] 
(8.5, -3.5) to[out=0, in=180] 
(11.4, -3.5);

\draw[-{angle 60}, black!30!white] (block_2.south east) to[out=-45, in=180] 
(6, -4.2) to[out=0, in=180] 
(12, -4.2) to[out=0, in=-135] 
(13.42, -3.8);

\node[gray] at (6.1,-3.9) {?};


%%% SW/FW and HW annotations

\draw[dashed] (10.5, 2.2) --++(0,-6.4);
\node[] at (7, 1.75) {Software \& firmware};
\node[] at (12, 1.75) {Hardware};


\end{tikzpicture}
}
    \end{figure}

    % \vspace{-2mm}
    A block is divided into warps. The warps of a block are associated with a single stream multiprocessor (SM) -- the GT 740 has 2 SMs, each with 192 cores.
    
\end{frame}


\begin{frame}{"Warping" can over provision GPU resources}
    Assume an input array with 384 elements and a GPU with 384 CUDA cores

    How many warps are spawned?
    \begin{itemize}\setlength\itemsep{0.5em}
        \item \inline{foo<<<2, 192>>>(...)} \uncover<2->{$\Rightarrow$ 12 (requires 384 cores)}
        \item \inline{foo<<<384, 1>>>(...)} \uncover<3->{$\Rightarrow$ 384 (requires 12288 cores)}
        \uncover<4->{\item \inline{foo<<<1, 192>>>(...)} $\Rightarrow$ 6 (requires 192 cores, 2 thread iterations)}
        \uncover<5->{\item \inline{foo<<<1, 384>>>(...)} $\Rightarrow$ 12 (requires 192 cores, 2 scheduling iterations)}
    \end{itemize}
    
\end{frame}


\begin{frame}[fragile]{Thread divergence within a warp}

    \begin{itemize}\setlength\itemsep{0.5em}
        \item Part of the warp has higher complexity, taking longer to execute.
        \item The other part of the warp becomes idle (has finished its job).
        \item Idle threads occupy cores and do not contribute to the processing of data.
    \end{itemize}

\vspace{2mm}

\begin{lstlisting}[language=C++]
__global__ void addOrMultiply(int *a, int *b, int *c) {
    int tid = blockIdx.x * blockDim.x + threadIdx.x;

    if (tid % 2 == 0) {     // If tid is even, add elements
        c[tid] = a[tid] + b[tid];
    } else {                // If tid is odd, multiply elements
        c[tid] = a[tid] * b[tid];
    }
}
\end{lstlisting}

\end{frame}


\begin{frame}{Striding}

    \begin{itemize}
        \item When the matching between threads and elements is not 1:1
        \item A thread may compute multiple elements
        \item \inline{if(abs_tid < input_size}{...} $\Rightarrow$ \inline{while(...){...}} or \inline{for(...){...}}
    \end{itemize}  

    \begin{figure}
        \centering
        \resizebox{0.75\textwidth}{!}{

\begin{tikzpicture}[]


\def\size{0.75}

\def\upperY{4}
\def\lowerY{0}

%%% DRAWING %%%%%%%%%%%%%%%%%%%%%%%%%%%%%%%%%%%%%%%%%%%%%%%%%%%%%%%%%%%%%%%%%


%%% Threads
\foreach \x in {0,...,3}{
    \draw[] (\x*\size, \upperY*\size) rectangle ++(\size, \size);
    % \node[anchor=center, align=center] at ((\x+0.5)*\size, 4.5*\size) {\x};
    \node[anchor=center, align=center] at (\x*\size+0.5*\size, \upperY*\size + 0.5*\size) {\x};
}


%%% Elements
\foreach \x in {0,...,9}{
    \draw[] (\x*\size, \lowerY) rectangle ++(\size, \size);
    \node[anchor=center, align=center] at (\x*\size+0.5*\size, \lowerY+0.5*\size) {\x};
}


%%% Lines
\foreach \x in {0,...,3}{
    \draw[-{angle 60}] (\x*\size+0.5*\size, \upperY*\size) to[out=-90, in=90] (\x*\size+0.5*\size, \lowerY*\size + \size);
    \draw[magenta, -{angle 60}] (\x*\size+0.5*\size, \upperY*\size) to[out=-90, in=90] (\x*\size+4.5*\size, \lowerY*\size + \size);
}
\foreach \x in {0,...,1}{
    \draw[blue, -{angle 60}] (\x*\size+0.5*\size, \upperY*\size) to[out=-90, in=90] (\x*\size+8.5*\size, \lowerY*\size + \size);
}


%%% Annotations
\node[anchor=east, align=right] at (-0.25, \upperY*\size+0.5*\size) {Grid of threads:};
\node[anchor=east, align=right] at (-0.25, \lowerY*\size+0.5*\size) {Input array:};

\node[black] (abs_tid) at (\size*6,\upperY*\size + 0.5*\size) {abs\_tid};
\draw[-{angle 60}](abs_tid.west) --++ (-\size*1.2, 0);
\node[black] (index) at (\size*12,\lowerY*\size + 0.5*\size) {index};
\draw[-{angle 60}](index.west) --++ (-\size*1.4, 0);

\node[black, anchor=west, align=left] at (\size*10, \upperY*\size + 0.5*\size) {Iteration 1 (k=0)};
\node[magenta, anchor=west, align=left] at (\size*10, \upperY*\size - 0.5*\size) {Iteration 2 (k=1)};
\node[blue, anchor=west, align=left] at (\size*10, \upperY*\size - 1.5*\size) {Iteration 3 (k=2)};


\end{tikzpicture}
}
    \end{figure}
    
    \inline{index = abs_tid + k * grid_size}, where \inline{index < input_size}
    
\end{frame}


\begin{frame}{Striding, divergence, and warps on an example}
    
    Assume we're using the GT 740 (2 SMs, each with 192 cores) and:
    \begin{itemize}\setlength\itemsep{0.5em}\vspace{1mm}
        \item We spawn a kernel with a single block of 256 threads
        \item To process an input array with 260 elements
        \item Striding is implemented
    \end{itemize}

    \vspace{3mm}
    \begin{itemize}\setlength\itemsep{0.5em}
        \uncover<2->{\item Is there thread divergence?}\uncover<3->{ {\color{Gray} Yes, at least one warp of 32 threads has to run twice, while only 4 threads are needed. Divergence can be avoided by using inputs with a size that is a multiple of 32.}}
        \uncover<4->{\item How many cores are utilized?}\uncover<5->{ {\color{Gray} Only 192 cores are utilized since a block (all of its warps) gets mapped to one SM.}}
    \end{itemize}
    
\end{frame}


\section{Exercise}

\begin{frame}{Goal}

    \textbf{Part 1 -- warps}
    \begin{itemize}\setlength\itemsep{0.5em}\vspace{1mm}
        \item Implement a kernel that inverts an image (calculate $x = 255 - x$, where $x \in (0, 255)$, for each color component -- RGB).
        \item Assess how the number of threads (full/partial warps) influences the computation time.
    \end{itemize}

    \vspace{4mm}

    \textbf{Part 2 -- thread divergence}
    \begin{itemize}\setlength\itemsep{0.5em}\vspace{1mm}
        \item Augment the inversion kernel, so that the R color component is calculated as $x = (x \, \% \, 25) \cdot 10$, $x \in (100, 200)$
        \item Assess how thread divergence impacts the computation time.
    \end{itemize}

    % {\footnotesize \textit{I encourage you to discuss all parts in the report, even if not fully implemented.}}
    
\end{frame}


\begin{frame}{Some pointers - nonobligatory}

\begin{itemize}\setlength\itemsep{1em}
    \item Implement the inverse kernel and use the provided image load/save function (Toledo $\rightarrow$ Docs) to invert an image. You can use as many threads as there are input elements.
    \item Implement striding in the above kernel $\rightarrow$ if difficult, first implement striding in a simpler kernel that handles a short 1D sequence, for example, using 12 threads, 16 elements (0, 1, \dots, 15), and adding 1 to each element.
    \item Check that the inversion kernel, with the added striding, generates a fully-inverted image even when using less threads than there are input data elements. For example, consider a ratio of 1:16 (threads:elements).
\end{itemize}
    
\end{frame}


\begin{frame}{Some pointers - nonobligatory}

\begin{itemize}\setlength\itemsep{1em}
    \item Invoke the inverse kernel from a loop, where you progressively use more threads. Time the kernel on every iteration (see cheat sheet).
    \item Get the average execution time, as the time may vary. For example, due to preemption by the OS. 
    \item Note that on some setups, the GPU may take more time to finish when first invoked -- you can avoid this by calling the function once before actually timing it.
\end{itemize}
    
\end{frame}


\begin{frame}{Some pointers - nonobligatory}

\begin{itemize}\setlength\itemsep{1em}
    \item Make a second kernel, where the R color component is processed according to the provided equation. 
    \item Check that the returned image is indeed different.
    \item Now make a second kernel, which instead of processing \inline{RGBRGB...RGB} operates on an input format of \inline{RR...RGGG...GBB...B} (re-format the data correspondingly).
    \item Evaluate how the execution time changes with / without thread divergence.
\end{itemize}
    
\end{frame}


\end{document}
\fancyarrow{Hello}{World}.